\documentclass[12pt, a4paper]{article}
\usepackage[utf8]{inputenc}
\usepackage[spanish]{babel}
\usepackage{graphicx}
\usepackage{float}
\usepackage{caption}
\usepackage{subcaption}
\usepackage{geometry}
\usepackage{listings}
\usepackage{xcolor}
\usepackage{amsmath}
\geometry{margin=2.5cm}

\lstset{
    language=Python,
    basicstyle=\ttfamily\small,
    keywordstyle=\color{blue},
    commentstyle=\color{gray},
    stringstyle=\color{orange},
    breaklines=true,
    frame=single
}

\title{Trabajo Práctico 1 \\ Filtros y Puerto Serie}
\author{Tu Nombre}
\date{\today}

\begin{document}
\maketitle
\tableofcontents
\newpage

% -------------------------------------------------------
\section{Introducción}

Este trabajo práctico integra dos temas vistos en el curso: el diseño y análisis
de filtros digitales en Python, y la comunicación por puerto serie. El objetivo
es controlar los parámetros de un filtro de coseno realzado (RC) y de raíz de
coseno realzado (RRC) mediante comandos enviados por puerto serie, y analizar
su comportamiento a través de distintas visualizaciones.

% -------------------------------------------------------
\section{Marco Teórico}

\subsection{Filtro de Coseno Realzado (RC)}

El filtro de coseno realzado es ampliamente utilizado en comunicaciones digitales
para limitar el ancho de banda de la señal transmitida minimizando la interferencia
entre símbolos (ISI, \textit{Inter-Symbol Interference}). Su respuesta en frecuencia
está definida por:

\begin{equation}
H(f) =
\begin{cases}
T, & |f| \leq \frac{1-\alpha}{2T} \\
\frac{T}{2}\left[1 + \cos\left(\frac{\pi T}{\alpha}\left(|f| - \frac{1-\alpha}{2T}\right)\right)\right], & \frac{1-\alpha}{2T} < |f| \leq \frac{1+\alpha}{2T} \\
0, & |f| > \frac{1+\alpha}{2T}
\end{cases}
\end{equation}

donde $T$ es el período de símbolo y $\alpha$ es el factor de roll-off
($0 \leq \alpha \leq 1$).

\subsection{Factor de Roll-off ($\alpha$)}

El factor de roll-off controla la transición entre la banda de paso y la banda
de rechazo del filtro. Sus efectos son:

\begin{itemize}
    \item $\alpha \to 0$: el filtro se aproxima a un filtro ideal de banda estrecha.
    La respuesta temporal presenta oscilaciones muy pronunciadas (alta sensibilidad
    al sincronismo) pero ocupa el mínimo ancho de banda posible.
    \item $\alpha = 1$: el filtro ocupa el doble del ancho de banda mínimo, pero
    la respuesta temporal decae mucho más rápido, siendo más tolerante a errores
    de sincronismo.
\end{itemize}

\subsection{Filtro de Raíz de Coseno Realzado (RRC)}

El filtro RRC es la raíz cuadrada espectral del filtro RC. Se utiliza en pares:
un filtro RRC en el transmisor y otro en el receptor. La convolución de ambos
produce un filtro RC completo, garantizando la condición de Nyquist para ISI nula.
A diferencia del RC, la respuesta al impulso del RRC \textbf{no} presenta cruces
por cero en los instantes de símbolo $t = \pm nT$.

\subsection{Parámetro span}

El parámetro span define la duración del filtro en cantidad de símbolos. Como el
filtro RC/RRC tiene soporte teóricamente infinito, en la práctica se trunca a
$\text{span} \times T$ segundos. Un span corto reduce el costo computacional pero
introduce distorsión por truncamiento, visible como ripple en la respuesta en
frecuencia. Un span largo aproxima mejor el comportamiento ideal.

\subsection{Parámetro sps (Samples Per Symbol)}

El parámetro sps define cuántas muestras se toman por período de símbolo $T$.
Determina la resolución de la representación discreta del filtro. Un valor de
sps bajo hace visibles las muestras individuales del filtro como una secuencia
discreta. Un valor alto produce una curva visualmente continua pero con mayor
cantidad de coeficientes.

% -------------------------------------------------------
\section{Implementación}

\subsection{Protocolo de comandos por puerto serie}

Se implementó un protocolo de comandos en texto plano sobre puerto serie
(simulado con \texttt{loop://} de pyserial). Los comandos disponibles son:

\begin{center}
\begin{tabular}{ll}
\hline
\textbf{Comando} & \textbf{Acción} \\
\hline
\texttt{set alpha <valor>} & Modifica el factor de roll-off (float) \\
\texttt{set span <valor>}  & Modifica la duración en símbolos (int) \\
\texttt{set sps <valor>}   & Modifica las muestras por símbolo (int) \\
\texttt{set rrc <True|False>} & Selecciona RRC o RC \\
\texttt{plot time}         & Respuesta al impulso (continua) \\
\texttt{plot freq}         & Respuesta en frecuencia \\
\texttt{plot both}         & Temporal y frecuencia en una figura \\
\texttt{plot discrete}     & Respuesta al impulso (discreta, stems) \\
\texttt{plot multi}        & Múltiples filtros con distintos $\alpha$ \\
\texttt{get coef}          & Imprime los coeficientes FIR \\
\texttt{status}            & Muestra parámetros actuales \\
\texttt{exit}              & Cierra el puerto y termina \\
\hline
\end{tabular}
\end{center}

\subsection{Descripción del script principal}

El script \texttt{tp\_main.py} instancia un objeto \texttt{RaisedCosineFilter}
con parámetros iniciales y lo reconfigura ante cada comando \texttt{set},
reconstruyendo el filtro con los nuevos valores. Las funciones de plot adicionales
(\texttt{plot\_both}, \texttt{plot\_discrete}, \texttt{plot\_multi}) fueron
implementadas fuera de la clase para no modificar \texttt{filtro.py}.

% -------------------------------------------------------
\section{Experimentos y Resultados}

\subsection{Efecto de sps con múltiples valores de $\alpha$}

\textbf{Motivación:} Se buscó analizar cómo la resolución de muestreo afecta
la representación del filtro para distintos roll-offs. Se fijaron span=6 y
rrc=True, y se varió sps $\in \{2, 4, 8, 12, 16\}$. Con sps=2 se está en el
límite de Nyquist: apenas dos muestras por símbolo, lo que produce una
representación muy tosca del filtro independientemente del valor de $\alpha$.
Al aumentar sps la curva se suaviza y los lóbulos laterales del filtro se
definen con mayor precisión. A partir de sps=8 los cambios son poco perceptibles
visualmente, lo que sugiere que ese valor es suficiente para representar
fielmente el filtro en este contexto.

\begin{figure}[H]
    \centering
    \includegraphics[width=0.9\textwidth]{multi_sps2.png}
    \caption{plot multi con sps=2, span=6, rrc=True. La baja resolución distorsiona
    la forma del filtro para todos los valores de $\alpha$.}
\end{figure}

\begin{figure}[H]
    \centering
    \includegraphics[width=0.9\textwidth]{multi_sps4.png}
    \caption{plot multi con sps=4, span=6, rrc=True. Se comienza a distinguir
    la forma característica del filtro.}
\end{figure}

\begin{figure}[H]
    \centering
    \includegraphics[width=0.9\textwidth]{multi_sps8.png}
    \caption{plot multi con sps=8, span=6, rrc=True. Representación típica,
    buen balance entre resolución y cantidad de coeficientes.}
\end{figure}

\begin{figure}[H]
    \centering
    \includegraphics[width=0.9\textwidth]{multi_sps12.png}
    \caption{plot multi con sps=12, span=6, rrc=True.}
\end{figure}

\begin{figure}[H]
    \centering
    \includegraphics[width=0.9\textwidth]{multi_sps16.png}
    \caption{plot multi con sps=16, span=6, rrc=True. La curva es prácticamente
    idéntica a sps=12, confirmando que aumentar sps más allá de 8-10 no aporta
    información nueva.}
\end{figure}

\subsection{Efecto de span con múltiples valores de $\alpha$}

\textbf{Motivación:} Se analizó el efecto del truncamiento del filtro. Se fijaron
sps=8 y rrc=True, y se varió span $\in \{2, 4, 6, 9, 12\}$. El filtro RC/RRC
tiene soporte infinito en el tiempo, por lo que al truncarlo se introduce un
error. Con span=2 el filtro es muy corto y la respuesta en frecuencia presenta
ripple significativo, alejándose del ideal. Al aumentar el span el filtro captura
más de la respuesta al impulso y la respuesta en frecuencia converge al
comportamiento teórico. La comparación entre span=2 y span=12 ilustra claramente
el compromiso entre longitud del filtro y fidelidad de la respuesta.

\begin{figure}[H]
    \centering
    \includegraphics[width=0.9\textwidth]{multi_span2.png}
    \caption{plot multi con span=2, sps=8, rrc=True. Truncamiento agresivo:
    la respuesta en frecuencia se aleja notablemente del ideal para todos los $\alpha$.}
\end{figure}

\begin{figure}[H]
    \centering
    \includegraphics[width=0.9\textwidth]{multi_span4.png}
    \caption{plot multi con span=4, sps=8, rrc=True.}
\end{figure}

\begin{figure}[H]
    \centering
    \includegraphics[width=0.9\textwidth]{multi_span6.png}
    \caption{plot multi con span=6, sps=8, rrc=True. Valor de referencia típico.}
\end{figure}

\begin{figure}[H]
    \centering
    \includegraphics[width=0.9\textwidth]{multi_span9.png}
    \caption{plot multi con span=9, sps=8, rrc=True.}
\end{figure}

\begin{figure}[H]
    \centering
    \includegraphics[width=0.9\textwidth]{multi_span12.png}
    \caption{plot multi con span=12, sps=8, rrc=True. La respuesta en frecuencia
    converge al comportamiento teórico del filtro RC/RRC.}
\end{figure}

\subsection{Comparación RRC vs RC con plot both}

\textbf{Motivación:} Se comparó el filtro RRC (rrc=True) contra el RC (rrc=False)
con los mismos parámetros ($\alpha=0.5$, span=6, sps=8). El RC presenta cruces
por cero exactos en los instantes de símbolo $t = \pm 1, \pm 2, \ldots$, lo
que garantiza ISI nula cuando se muestrea en el instante correcto. El RRC en
cambio no presenta esa propiedad por sí solo: necesita ser convolucionado con
otro RRC (en el receptor) para recuperar el RC completo. En frecuencia, ambos
tienen el mismo soporte espectral, pero la magnitud del RRC es la raíz cuadrada
de la del RC.

\begin{figure}[H]
    \centering
    \includegraphics[width=0.9\textwidth]{both_rrc.png}
    \caption{plot both con rrc=True ($\alpha=0.5$, span=6, sps=8). Filtro RRC:
    respuesta temporal sin cruces por cero en instantes de símbolo.}
\end{figure}

\begin{figure}[H]
    \centering
    \includegraphics[width=0.9\textwidth]{both_rc.png}
    \caption{plot both con rrc=False ($\alpha=0.5$, span=6, sps=8). Filtro RC:
    se observan los cruces por cero en $t = \pm 1, \pm 2, \ldots$}
\end{figure}

\subsection{Respuesta discreta con plot discrete}

\textbf{Motivación:} El filtro FIR es en esencia una secuencia finita de
coeficientes (muestras discretas). Con sps alto esta naturaleza discreta queda
oculta porque las muestras se ven como una curva continua. Al reducir sps=4
cada muestra es claramente visible como un valor individual, recordando que
lo que se implementa en hardware o software es un vector de números, no una
función continua. Se usó $\alpha=0.5$, span=6, rrc=True.

\begin{figure}[H]
    \centering
    \begin{subfigure}{0.48\textwidth}
        \includegraphics[width=\textwidth]{discrete_sps4.png}
        \caption{sps=4: cada muestra es visible individualmente}
    \end{subfigure}
    \hfill
    \begin{subfigure}{0.48\textwidth}
        \includegraphics[width=\textwidth]{discrete_sps8.png}
        \caption{sps=8: los stems quedan más juntos, la curva parece continua}
    \end{subfigure}
    \caption{plot discrete con $\alpha=0.5$, span=6, rrc=True. La comparación
    entre sps=4 y sps=8 muestra cómo sps controla la densidad de muestras
    del filtro FIR.}
\end{figure}

% -------------------------------------------------------
\section{Conclusiones}

A partir de los experimentos realizados se puede concluir que:

\begin{itemize}
    \item El factor $\alpha$ controla el compromiso entre ancho de banda ocupado
    y sensibilidad al sincronismo. Valores bajos son más eficientes espectralmente
    pero más exigentes en precisión de temporización.
    \item El parámetro span determina la calidad de la aproximación al filtro ideal.
    Un span reducido introduce ripple en frecuencia; valores a partir de 6 suelen
    ser suficientes en la práctica.
    \item El parámetro sps no altera la forma del filtro sino su resolución de
    representación. A partir de sps=8 los cambios visuales son mínimos.
    \item El filtro RRC no cumple por sí solo la condición de Nyquist; debe
    usarse en par (transmisor y receptor) para obtener el RC completo y garantizar
    ISI nula.
    \item El plot discrete hace evidente que un filtro FIR es una secuencia finita
    de coeficientes, concepto clave para su implementación en hardware o software.
\end{itemize}

\end{document}